\chapter{Análise de Dados e Automação com LLMs}
\label{ch:analise-dados}

\ficcao O relatório de vendas da Lavandaria mostra 90 euros. No caderno estão dois pedidos de 30 euros. A consulta correu sem erros. Antes de pedir ao modelo uma explicação para o aumento das vendas, vale a pena perceber por que motivo um pedido entrou duas vezes na soma.

\section{Saber o que representa uma linha}

Há uma linha por pedido na primeira tabela e uma por artigo na segunda. O primeiro pedido tem dois artigos. Se juntares as tabelas, o total do pedido aparece uma vez em cada linha de artigo. O cálculo pode estar certo para a tabela produzida e errado para a pergunta original.

Podes fazer a conta sem computador. Depois da junção, vês 30, 30 e 30. A soma dá 90. Eliminar valores repetidos dá 30, porque os dois pedidos legítimos têm o mesmo total. A pergunta exige contar cada pedido uma vez, usando o identificador do pedido para distinguir os dois.

\section{Executar a consulta num conjunto pequeno}

SQL é uma linguagem para pedir operações a uma base de dados. Podes repetir o exemplo no SQLite, um sistema que guarda as tabelas localmente. Cria-as e começa pela consulta que as junta. Dentro de \texttt{SUM}, o operador \texttt{DISTINCT} elimina valores iguais antes da soma~\citep{sqliteagregados2026}.

\begin{lstlisting}
CREATE TABLE pedidos (id INTEGER PRIMARY KEY, total INTEGER);
CREATE TABLE artigos (pedido_id INTEGER, nome TEXT);
INSERT INTO pedidos VALUES (1, 30), (2, 30);
INSERT INTO artigos VALUES (1, 'meia'), (1, 'saco'),
                           (2, 'toalha');

SELECT COUNT(*) AS linhas, SUM(p.total) AS soma,
       SUM(DISTINCT p.total) AS soma_distinta
FROM pedidos AS p
JOIN artigos AS a ON a.pedido_id = p.id;
\end{lstlisting}

\noindent\begin{minipage}{\linewidth}
A segunda consulta usa \texttt{EXISTS} para verificar se o pedido tem pelo menos um artigo. A soma continua a usar uma linha por pedido.
\begin{lstlisting}
SELECT COUNT(*) AS pedidos, SUM(p.total) AS soma
FROM pedidos AS p
WHERE EXISTS (
    SELECT 1 FROM artigos AS a WHERE a.pedido_id = p.id
);
\end{lstlisting}
\end{minipage}

Na execução local, a primeira consulta devolveu $(3,90,30)$ e a segunda $(2,60)$. Para incluir pedidos sem artigos, consulta só a tabela de pedidos. Se também precisares de estatísticas de artigos, agrega primeiro por pedido e verifica a cardinalidade antes de juntar.

\begin{promptlabbox}{Rever uma consulta com um resultado conhecido}
\textbf{Objetivo e contexto:} Verificar linhas e totais antes de aplicar uma consulta a dados maiores. Usa o esquema fictício acima ou uma base de teste autorizada.

Escreve a resposta à mão. Pede SQL para o dialeto que vais usar e indica o que representa uma linha de cada tabela. Compara com a tua consulta. Acrescenta um pedido sem artigos e decide se deve entrar no total. Se a pergunta usar datas, testa os limites inicial e final do intervalo.

\textbf{Resultado esperado e validação:} A consulta devolve o número de pedidos e o total definidos pela regra. Explica o efeito de uma duplicação, de um nulo e de uma exclusão. Confirma permissões, volume e tempo de execução antes de passar a outro ambiente.

\textbf{Falhas possíveis:} Comparar apenas o número de linhas, deduplicar valores quando a unidade é o pedido ou executar uma consulta sugerida sem a rever.
\end{promptlabbox}

\section{Explorar antes de escrever a interpretação}

Começa por contagens, ausências, intervalos, unidades e grupos. Se duas colunas se moverem em conjunto, considera explicações alternativas e escolhe a análise que as separaria. Um modelo pode ajudar a propor perguntas e consultas. Guarda as tabelas e os gráficos que efetivamente calculaste para conseguir regressar ao resultado.

Para períodos consecutivos, um intervalo com início incluído e fim excluído permite pôr um registo de fronteira num só período. Define o fuso horário se houver instantes de tempo. Regista como trataste cancelamentos, dados ausentes e moedas. Um valor desconhecido tem uma consequência diferente de um valor zero, sobretudo quando calculas médias.

\section{Dar ao relatório métricas com contexto}

Um relatório pode receber um objeto JSON com período, unidade, numerador e denominador. Este exemplo sintético distingue o universo da amostra que respondeu.

\begin{lstlisting}
{
  "periodo": "2026-10",
  "saidas_registadas": 20,
  "respostas_ao_inquerito": 8,
  "respostas_que_mencionam_preco": 5
}
\end{lstlisting}

«Cinco das oito respostas mencionam preço» descreve os dados. Para dizer o que aconteceu nas vinte saídas, faltam doze respostas. Para atribuir a causa de uma saída, seria preciso um desenho que sustentasse essa conclusão. Pede que cada número do relatório aponte ao campo ou cálculo correspondente.

\begin{promptlabbox}{Ligar dados e narrativa}
\textbf{Objetivo e contexto:} Produzir um relatório curto a partir de métricas calculadas, com unidades e denominadores visíveis.

Usa o JSON acima. Redige três frases que distingam contagem, proporção e dúvida. Se recorreres a um modelo, confronta o rascunho com uma versão feita à mão. Testa também o caso de zero respostas ao inquérito.

\textbf{Resultado esperado e validação:} A proporção entre respondentes é $5/8=62{,}5\%$. Com zero respostas, essa proporção fica indefinida. Cada frase indica o universo a que se aplica.

\textbf{Falhas possíveis:} Dividir pelo universo errado, trocar pontos percentuais por variação relativa ou escrever uma causa que os dados não mediram.
\end{promptlabbox}

\section{Automatizar a parte que tem uma regra}

Para normalizar cabeçalhos de CSV, escreve regras de caixa, acentos e separadores e testa colisões. «Preço total» e «Preco total» podem produzir o mesmo identificador. Decide como preservar a relação com os nomes de origem. Para aproximar categorias entre sistemas, pede uma proposta de correspondência com a opção «sem correspondência». Guarda a aprovação e versiona o mapa antes de o aplicar a outros registos. O capítulo~\ref{ch:software} mostra como transformar essas regras em casos verificáveis.
